\documentclass[10pt,twocolumn]{article}
\usepackage[utf8]{inputenc}\usepackage[T1]{fontenc}\usepackage[english]{babel}\usepackage{amsmath,amssymb}\usepackage[margin=1.8cm]{geometry}
\usepackage{times}\usepackage{hyperref}
\title{\textbf{Exact two-sided certificates for maximal quantum Bell violations: an agentic laboratory study of the doubly tilted CHSH family}}
\author{Team Ninja Turtles\\ \small Hack-Nation 2026, Challenge 3 (Agentic Scientific Discovery)}
\date{Preprint, \today}
\begin{document}\maketitle
\section{Abstract}

We study the doubly tilted CHSH family with correlator coefficients a = b = [p, 0] and c = [[1, 1], [1, -1]]. An automated laboratory computed lower bounds on the maximal quantum value Q(p) from explicit strategies and upper bounds from the NPA hierarchy. A verifier certified both bounds in exact rational arithmetic. Q is certified to within 1e-6 at p = 0.3 {\scriptsize[C-bell-R4]}, at p = 1/2 {\scriptsize[C-bell-R1]} and at p = 131/200 {\scriptsize[C-bell-R6]}. At p = 0.95 and p = 0.99 the explicit strategies violate the classical bound L = 2+2p, but the two-sided gap is not closed {\scriptsize[C-bell-R2, C-bell-R2-RT2]}. At p = 131/200, NPA level 1+AB is strictly not tight, while level 2 is {\scriptsize[C-kombi1]}. No closed formula for Q(p) was established. The preregistered anchor hypothesis H6a failed {\scriptsize[C-H6a]}, and no novelty is claimed {\scriptsize[C-H6b]}.

\section{1 Introduction}

\textbf{Question.} How does the maximal quantum value Q(p) behave for the doubly tilted CHSH family with a = b = [p, 0]? Do qubits suffice? Where does the NPA level 1+AB cease to be tight? {\scriptsize[C-bell-R1, C-bell-R5]}

\textbf{Contribution.} The study reports individually certified values and gap bounds. It also lists claims that could not be established.

\textbf{Results.}
\begin{itemize}
\item Two-sided certificates for Q at three parameter values: p = 0.3 {\scriptsize[C-bell-R4]}, p = 1/2 {\scriptsize[C-bell-R1]} and p = 131/200 {\scriptsize[C-bell-R6]}.
\item A rigorous quantum violation of L at p = 0.95 and p = 0.99, without a closed two-sided gap {\scriptsize[C-bell-R2, C-bell-R2-RT2]}.
\item Strict non-tightness of NPA level 1+AB at p = 131/200 {\scriptsize[C-kombi1]}.
\item Verifier lower bounds on the gap between the 1+AB optimum and the best certified strategy at six parameter values {\scriptsize[C-kombi2]}.
\item A list of negative results, red-team outcomes and unmet preregistered hypotheses, given in the section on negative results {\scriptsize[C-neg1, C-H6a]}.
\end{itemize}

\section{2 Setting and certificates}

The Bell inequality is given in correlator form with coefficients a = [p, 0], b = [p, 0], c = [[1, 1], [1, -1]] and constant 0 {\scriptsize[C-bell-R1]}. The classical bound is L = 2+2p, and at p = 1/2 it is L = 3 {\scriptsize[C-bell-R2]}.

Lower bounds on Q come from explicit quantum strategies of dimension d per side, found by see-saw optimisation. The verifier evaluates their Bell value exactly in rational arithmetic {\scriptsize[C-bell-R1-RT1]}. Upper bounds come from NPA levels 1+AB, 2 and 2+AAB, certified by rational dual certificates {\scriptsize[C-bell-R1, C-bell-R2-RT1]}. Q counts as proved when the gap between lower and upper bound is below 1e-6 {\scriptsize[C-bell-R1]}. Rational feasible moment matrices give lower bounds on the NPA optimum, and the verifier compares these with the best exact strategy {\scriptsize[C-bell-R5]}.

\section{3 Method: the agentic laboratory}

The laboratory ran 8 rounds, with 7 checked statements and 1 negative result, at a cost of 4.31 USD {\scriptsize[C-methode]}. Each round was preregistered before the experiment {\scriptsize[C-methode]}. Researcher agents proposed claims, and a code verifier certified or rejected them. A red team attacked each claim with counter-checks. Agent interpretations are labelled as unverified hypotheses and are not used here {\scriptsize[C-bell-R1-I]}.

Before the run, the verifier passed a self-test with 22/22 cases correct (10 true, 12 false statements or rule violations) {\scriptsize[C-selbsttest]}. The anchors were CHSH (classical bound 2, quantum value 2.828427125), the chain with 4 settings (classical bound 6), I3322 (classical bound 0, upper bound 0.2509 at NPA level 2+AAB) and tilted CHSH with formula sqrt(8+2\emph{p}*2) {\scriptsize[C-selbsttest]}.

\section{4 Results}

\textbf{Proposition A (computed_rigorous).} For p = 1/2, Q = 3.13693731 up to 1e-6 {\scriptsize[C-bell-R1]}. The verifier certified 3.13693731130907 ≤ Q ≤ 3.13693732592937 (gap 1.46e-08; strategy d=2, NPA 1+AB) {\scriptsize[C-bell-R1, C-bell-R3]}. The level 2+AAB upper bound is 3.13693731768845 {\scriptsize[C-bell-R3-RT2]}. An independent see-saw strategy gave the exact rational Bell value 3.13693731130907 with d=2 {\scriptsize[C-bell-R1-RT1]}.

\textbf{Proposition B (computed_rigorous).} For p = 0.3, Q = 2.93828832 up to 1e-6, with 2.93828833175884 ≤ Q ≤ 2.938288368081 (gap 3.63e-08; strategy d=2, NPA 1+AB) {\scriptsize[C-bell-R4]}. An independent d=2 strategy gave the exact value 2.93828833175884 {\scriptsize[C-bell-R4-RT2]}.

\textbf{Proposition C (computed_rigorous).} For p = 131/200, Q = 3.35825951388 up to 1e-6, with 3.35825951388355 ≤ Q ≤ 3.35825955532047 (gap 4.14e-08; strategy d=2, NPA 2) {\scriptsize[C-bell-R6]}. An independent d=2 strategy gave the exact value 3.35825951388355 {\scriptsize[C-bell-R6-RT2]}.

\textbf{Proposition D (computed_rigorous).} The classical bound is exceeded at two parameter values.
\begin{itemize}
\item At p = 0.99, the verifier's explicit strategy has exact value 3.98000132884628, and L = 199/50 {\scriptsize[C-bell-R2]}. The level 2+AAB upper bound is 3.98000623486376 {\scriptsize[C-bell-R2-RT1]}.
\item At p = 0.95, L = 3.9 {\scriptsize[C-bell-R2-RT2]}. Also 3.90016389936512 ≤ Q ≤ 3.90019981472908 (gap 3.59e-05) {\scriptsize[C-bell-R2-RT2]}. The gap is too large for certification of Q, but the lower bound exceeds L {\scriptsize[C-bell-R2-RT2]}.
\item At p = 1, L = 4 and the verifier's quantum strategy attains exactly 4, so no violation was found there {\scriptsize[C-bell-R1-RT2]}.
\end{itemize}

\textbf{Proposition E (computed_rigorous).} At p = 131/200, Q ≤ 3.35825955532047 (NPA level 2, dual certificate), while the optimum of NPA level 1+AB is ≥ 3.35826280973 (rational feasible moment matrix) {\scriptsize[C-kombi1]}. Level 1+AB is therefore strictly not tight there: its optimum exceeds Q by at least 3.25e-06 {\scriptsize[C-kombi1]}.

\textbf{Proposition F (computed_rigorous).} Lower bounds on (NPA 1+AB optimum) minus (best exactly certified strategy) are as follows {\scriptsize[C-kombi2]}:

| p | lower bound on gap | source |
|---|---|---|
| 0.5 | -1.726e-08 | {\scriptsize[C-kombi2]} |
| 0.62 | -2.336e-08 | {\scriptsize[C-kombi2]} |
| 0.652 | 1.843e-07 | {\scriptsize[C-kombi2]} |
| 0.6532 | 1.025e-06 | {\scriptsize[C-kombi2]} |
| 0.6533 | 1.120e-06 | {\scriptsize[C-kombi2]} |
| 0.655 | 3.296e-06 | {\scriptsize[C-kombi2]} |

Up to p = 0.62 the bound is below 1e-7 in magnitude, which is the accuracy of the SDP solvers {\scriptsize[C-kombi2]}. From p = 0.652 on it is positive and grows monotonically over the points examined {\scriptsize[C-kombi2]}. A threshold defined by a fixed tolerance such as 1e-6 therefore depends on that tolerance {\scriptsize[C-kombi2]}. At p = 0.62, the strategy value 3.30352972153 (d≤3) and the 1+AB value ≥ 3.30352969818 agree to within solver accuracy {\scriptsize[C-bell-R5-RT1]}. At p = 0.72, NPA level 2 reaches ≥ 3.46649094809 against a best strategy of 3.46649100491 (d≤2) {\scriptsize[C-bell-R5-RT2]}.

\textbf{Summary of certified values.}

| p | Lower bound (strategy) | Upper bound (NPA) | Gap | Status | Source |
|---|---|---|---|---|---|
| 0.3 | 2.93828833175884 (d=2) | 2.938288368081 (1+AB) | 3.63e-08 | proved | {\scriptsize[C-bell-R4]} |
| 1/2 | 3.13693731130907 (d=2) | 3.13693732592937 (1+AB) | 1.46e-08 | proved | {\scriptsize[C-bell-R1]} |
| 131/200 | 3.35825951388355 (d=2) | 3.35825955532047 (2) | 4.14e-08 | proved | {\scriptsize[C-bell-R6]} |
| 0.95 | 3.90016389936512 (d=2) | 3.90019981472908 (2+AAB) | 3.59e-05 | not certified | {\scriptsize[C-bell-R2-RT2]} |
| 0.99 | 3.98000132884628 | 3.98000623486376 (2+AAB) | not stated | not certified | {\scriptsize[C-bell-R2, C-bell-R2-RT1]} |

\section{5 Negative results, preregistration and red team}

\textbf{Negative result.} The scan for the transition at which NPA 1+AB ceases to be tight, over the points p = 0.6, 0.63, 0.65, 0.655, 0.66, 0.7, 0.8, 0.9 and 0.99, produced no claim that passed verification {\scriptsize[C-neg1]}.

\textbf{Preregistration.} H6a (at least 2 of 3 anchor questions confirmed two-sided) was not met: the integrator chose 0 of the 3 anchor questions, and 0 confirmed statements concern anchor questions {\scriptsize[C-H6a]}. H6b (at least one certified statement about an inequality outside the named families) was met, with 6 confirmed statements concerning the doubly tilted CHSH family or other unnamed inequalities {\scriptsize[C-H6b]}. Novelty relative to the literature is not claimed {\scriptsize[C-H6b]}.

\textbf{Red team.} Red-team checks targeting the confirmed values failed in the sense that the attacks did not succeed.
\begin{itemize}
\item No strategy exceeded the certified values at p = 1/2, p = 0.3 and p = 131/200 {\scriptsize[C-bell-R1-RT1, C-bell-R4-RT2, C-bell-R6, C-bell-R6-RT2]}.
\item The upper bounds were not undercut by any strategy {\scriptsize[C-bell-R3-RT2, C-bell-R2-RT1]}.
\item No violation of L was found at p = 1 {\scriptsize[C-bell-R1-RT2]}.
\item NPA 1+AB was not found loose at p = 1/2 or p = 0.62, and NPA level 2 was not found loose at p = 131/200 or p = 0.72 {\scriptsize[C-bell-R4-RT1, C-bell-R5-RT1, C-bell-R6, C-bell-R6-RT1, C-bell-R5-RT2]}.
\end{itemize}

One claim is contested {\scriptsize[C-bell-R8]}. At p = 0.6533, the 1+AB optimum lies at least 5e-07 above the best strategy found by the verifier in dimension up to 4 {\scriptsize[C-bell-R8]}. Here the verifier reports 3.35554254508 against 3.35554142523, a gap ≥ 1.120e-06 {\scriptsize[C-bell-R8]}. This does not show that Q itself lies below the 1+AB value {\scriptsize[C-bell-R8]}. A red-team check at p = 0.6532 passed (statement contested), with a gap ≥ 1.025e-06 {\scriptsize[C-bell-R8-RT1]}. A check at p = 0.652 did not show a gap of 1e-6, with a gap ≥ 1.843e-07 {\scriptsize[C-bell-R8-RT2]}. The agent's unverified interval for the transition is not used here.

\section{6 Limitations and open questions}

\begin{itemize}
\item \textbf{Not shown:} a closed formula for Q(p). The fit attempts addressed this question {\scriptsize[C-bell-R3, C-bell-R4]}, but no certified formula resulted.
\item \textbf{Not shown:} that Q > L for all p < 1, or that Q = L at p = 1. The claims certify a strategy value of 4 at p = 1 and violations only at the tested points {\scriptsize[C-bell-R1-RT2, C-bell-R2]}.
\item \textbf{Not shown:} that qubits suffice in general. Certified strategies at the tested points have d=2 or d≤3, and the lower bounds there match the upper bounds only to the stated gaps {\scriptsize[C-bell-R5-RT1, C-bell-R6]}.
\item A sharp transition value for NPA 1+AB is not established. The verifier bounds are tolerance-dependent {\scriptsize[C-kombi2]}, and the contested claim {\scriptsize[C-bell-R8]} is not resolved.
\item Q is not certified at p = 0.95 and p = 0.99 {\scriptsize[C-bell-R2-RT2, C-bell-R2-RT1]}.
\item No novelty is claimed beyond what the claims support {\scriptsize[C-H6b]}.
\end{itemize}

\textbf{Related work.} The NPA hierarchy yields exact bounds via an infinite series of semidefinite programs {\scriptsize[C-lit6]}. In the simplest Bell scenario, the correlations at levels 1+AB and 2 coincide under a plausible criterion {\scriptsize[C-lit13]}. Analytic criteria for the boundary of quantum correlations are hard to establish even there {\scriptsize[C-lit14]}. An iterative algorithm computes the quantum maximum of two-outcome Bell inequalities in a given Hilbert space dimension {\scriptsize[C-lit10]}. Operations on CHSH-type inequalities that preserve the Tsirelson bound but change the classical bound have been studied {\scriptsize[C-lit5]}. For I3322, finite-dimensional systems are conjectured not to attain the true quantum maximum {\scriptsize[C-lit9]}. Proposition E shows that, for a functional with marginal terms in the simplest scenario, the optima of NPA levels 1+AB and 2 differ {\scriptsize[C-kombi1]}. The coincidence of levels 1+AB and 2 reported in {\scriptsize[C-lit13]} is stated for correlations under a plausible criterion; whether its scope covers functionals with marginal terms was not checked beyond the abstract, so we do not claim a contradiction. Beyond this, we make no claim about how these results relate to the present family.

---
\emph{Verification log (automatic):} 34 claim IDs cited; unsupported statements per writing round: 30 → 2; 0 sentences removed; remaining violations: 0. Claim texts and evidence: `paper_belege.json`.

\end{document}
